\chapter{La tierra que desapareció}

\label{cap:tierra}
El capítulo \ref{cap:fincas} mostró de dónde vienen las unidades territoriales del departamento.
Éste mide qué quedó de ellas. Entre los censos agropecuarios de 2002 y 2018, el
departamento La Caldera perdió el ochenta y ocho por ciento de su superficie agropecuaria
mientras la provincia ganaba. Ese dato es el punto de partida de todo lo que el libro documenta después, porque parte de la tierra que sale del registro agropecuario es la que entra al mercado de suelo urbano; el apartado sobre el censo de 2008 precisa qué era esa superficie y cuánto de ella era monte que simplemente dejó de declararse.

\section{La estructura agraria medida}

El artículo 6º del proyecto de ordenanza enumera las fincas del municipio. El Censo Nacional Agropecuario 2018 permite medir lo que queda de esa estructura.

\begin{ficha}{Explotaciones agropecuarias, departamento La Caldera, CNA 2018}
\begin{itemize}[leftmargin=*,nosep]
\item \textbf{Total de explotaciones:} 131. De ellas, 70 con límites definidos o mixtas, y \textbf{61 sin límites definidos}.
\item \textbf{Superficie total de las explotaciones con límites:} 11.417,5 hectáreas, distribuidas en 78 parcelas. Se registran además 64 terrenos sin límites.
\item \textbf{Superficie departamental de referencia:} 1.091,9 km², unas 109.200 hectáreas (INDEC, Censo 2022; capítulo \ref{cap:fincas}). Las explotaciones agropecuarias cubren, por lo tanto, cerca del \textbf{10,5 por ciento} del departamento.
\end{itemize}
\textbf{Distribución por escala de extensión:}
\begin{center}\small
\begin{tabular}{@{}lrrr@{}}
\toprule
Escala & EAPs & Superficie (ha) & \% de la superficie \\
\midrule
Hasta 5 ha & 25 & 81,3 & 0,7 \\
5,1 a 10 ha & 12 & 87,0 & 0,8 \\
10,1 a 25 ha & 13 & 215,2 & 1,9 \\
25,1 a 50 ha & 8 & 287,0 & 2,5 \\
100,1 a 200 ha & \textit{s} & 320,0 & 2,8 \\
200,1 a 500 ha & 5 & 1.623,0 & 14,2 \\
500,1 a 1.000 ha & 4 & 2.934,0 & 25,7 \\
\textbf{5.000,1 a 7.500 ha} & \textit{s} & \textbf{5.870,0} & \textbf{51,4} \\
\midrule
\textbf{Total} & \textbf{70} & \textbf{11.417,5} & \textbf{100} \\
\bottomrule
\end{tabular}
\end{center}
\textit{La letra} s \textit{indica dato reservado por secreto estadístico: hay explotaciones en esa escala, pero son tan pocas que publicar el número las identificaría.}

\textit{Fuentes: la distribución por tipo de delimitación, de INDEC, \textit{Censo Nacional Agropecuario 2018. Resultados definitivos}, 2021, cuadro 2.1.16, pág.\ 100: 67 explotaciones con límites definidos (74 parcelas, 11.075,5 ha), 3 mixtas (4 parcelas, 342,0 ha y 3 terrenos sin límites) y 61 sin límites definidos; la escala de extensión, del sistema de consultas del mismo censo (cuadros C2 1 y C2 2). Cotejadas en septiembre de 2026.}
\end{ficha}

Una sola explotación agropecuaria concentra 5.870 hectáreas: el \textbf{51,4 por ciento de toda la superficie agropecuaria del departamento}. La escala está enmascarada por secreto estadístico, pero los totales del cuadro permiten deducir que en ese tramo hay una sola explotación: la misma que, más abajo, aparece como la única sobreviviente de los tramos mayores a mil hectáreas.

Los tres tramos superiores reúnen diez explotaciones y 10.427 hectáreas, el 91,3 por ciento. Las cincuenta y ocho explotaciones de los cuatro tramos inferiores --- todas las de menos de cincuenta hectáreas --- reúnen 670,5 hectáreas, el 5,9 por ciento.

Ésta es la estructura de tenencia que el capítulo describía desde el topónimo de Vaqueros y desde la lista de fincas del artículo 6º, ahora con números. La continuidad entre la merced colonial y el presente no es una figura retórica: la mitad del suelo agropecuario del departamento sigue en una sola mano.

El segundo dato es de otro orden y toca un hueco declarado de este libro.

\textbf{Sesenta y una de las 131 explotaciones del departamento --- el 46,6 por ciento --- no tienen límites definidos}, contra el 32,8 por ciento provincial. Se registran 64 terrenos sin límites.

La categoría censal ``sin límites definidos'' es técnica y no equivale a ``sin título''. El propio censo la define: son explotaciones cuyas parcelas no están delimitadas porque sus tierras forman parte de una \textbf{«unidad mayor»}, que puede ser un campo comunero, un campo perteneciente a comunidades de pueblos originarios, un parque o una reserva nacional, otro tipo de tierras fiscales o tierras privadas (INDEC, \textit{CNA 2018. Resultados definitivos}, nota 3, pág.\ 29). Es el indicador más próximo que este relevamiento encontró a la pregunta que el capítulo \ref{cap:resistencias} deja abierta sobre la existencia de tenencia precaria o posesión sin título en el departamento.

El censo no publica, por departamento, a qué tipo de dominio pertenecen esos terrenos. Sí lo publica por provincia, y la distribución salteña orienta la lectura: de los 4.392 terrenos sin límites definidos de Salta, 2.375 ---el 54 por ciento--- están en tierras de comunidades de pueblos originarios; 1.235 en tierras privadas; 485 en campos comunales o comuneros; 267 en otras tierras fiscales, y 30 en parques o reservas. Por régimen de tenencia, 1.817 corresponden a integrantes de una comunidad y 1.157 a ocupación de hecho (cuadro 3.3, págs.\ 145 a 147).

\begin{cautela}
\textbf{Una hipótesis que estos datos permiten formular, y sus límites.} La comunidad kolla Kondorwaira reclama como territorio, en su carpeta técnica, la finca Potrero de Castilla o Las Nieves, \textbf{catastro 102, de propiedad fiscal}; la integran treinta y una familias que crían ganado menor y mayor y cultivan papas, maíz y otros productos andinos, trasladándose entre la «casa grande» de verano y el «puesto» de invierno (Pérez Domínguez y Belmonte, «Arriba están los cerros y abajo la ciudad», \textit{Huellas}, 2024, \url{http://dx.doi.org/10.19137/huellas-2024-2819}; capítulo \ref{cap:resistencias}). Es exactamente la forma de producción que el censo registra como explotación sin límites definidos dentro de una unidad mayor fiscal o comunitaria.

\textbf{Este libro no puede afirmar} que las familias de la comunidad sean una parte de las sesenta y una explotaciones sin límites del departamento: el censo no publica el dominio por departamento y la cifra de explotaciones no tiene por qué coincidir con la de familias. Lo que sí puede afirmar es que la categoría existe para describir situaciones como ésa, y que en el departamento hay una. La verificación es un procesamiento especial del censo, pedido en el apéndice \ref{cap:pedidos}.
\end{cautela}

Casi la mitad de las unidades productivas caldereñas opera sobre tierra sin límites establecidos. Esa población no aparece en ninguna cobertura periodística ni en ningún expediente judicial relevado; la literatura académica sólo la alcanza, en parte y sin nombrarla así, en los trabajos sobre la comunidad Kondorwaira. El Censo Agropecuario la cuenta.

\section{La desaparición de la tierra agropecuaria}

El Censo Nacional Agropecuario 2002 publicó resultados por departamento. Con ellos, el capítulo deja de describir una estructura y pasa a medir un proceso.

\begin{ficha}{Departamento La Caldera, CNA 2002 y CNA 2018}
\begin{center}\small
\begin{tabular}{@{}lrrr@{}}
\toprule
 & 2002 & 2018 & Variación \\
\midrule
Explotaciones agropecuarias & 308 & 131 & $-57,5\%$ \\
\quad con límites definidos & 126 & 70 & $-44,4\%$ \\
\quad sin límites definidos & 182 & 61 & $-66,5\%$ \\
\textbf{Superficie (ha)} & \textbf{98.774,9} & \textbf{11.417,5} & $\mathbf{-88,4\%}$ \\
\bottomrule
\end{tabular}
\end{center}
\textit{Fuente: INDEC, Censos Nacionales Agropecuarios 2002 (cuadros 1.1 y 1.2, Salta) y 2018.}
\end{ficha}

\textbf{El departamento perdió 87.357 hectáreas de superficie agropecuaria en dieciséis años.} Sobre la superficie departamental de referencia, unas 109.200 hectáreas, las explotaciones pasaron de cubrir el \textbf{90 por ciento} del territorio al \textbf{10,5}. (Con las 146.300 hectáreas que este capítulo usaba antes, las proporciones eran del 67,5 y del 7,8; la caída relativa es la misma, porque no depende de la superficie del departamento.)

Una caída de esa magnitud obliga a preguntarse si el dato es real o si responde a un cambio de cobertura censal. El control es la propia provincia, y despeja la duda.

Entre 2002 y 2018 la superficie agropecuaria de Salta \textbf{creció} un 2,8 por ciento, de 4.269.499 a 4.387.088 hectáreas, mientras su cantidad de explotaciones caía un 15,5 por ciento. Es decir: en la provincia hubo concentración --- menos explotaciones sobre más superficie --- que es el proceso agrario argentino conocido.

En La Caldera no hubo concentración. Hubo desaparición. Las explotaciones cayeron casi cuatro veces más rápido que en la provincia, y la superficie no se concentró: se fue. La participación del departamento en la superficie agropecuaria provincial pasó del 2,31 por ciento al 0,26.

\subsection*{Dónde se fue la tierra}

La comparación por escala de extensión localiza la pérdida con precisión.

{\small
\begin{longtable}{@{}lrrrr@{}}
\toprule
Escala & EAP 2002 & ha 2002 & EAP 2018 & ha 2018 \\
\midrule
\endhead
Hasta 5 ha & 18 & 45,8 & 25 & 81,3 \\
5,1 a 10 & 4 & 33,0 & 12 & 87,0 \\
10,1 a 25 & 16 & 296,0 & 13 & 215,2 \\
25,1 a 50 & 14 & 502,5 & 8 & 287,0 \\
50,1 a 100 & 18 & 1.334,0 & 0 & --- \\
100,1 a 200 & 17 & 2.436,6 & \textit{s} & 320,0 \\
200,1 a 500 & 13 & 4.465,0 & 5 & 1.623,0 \\
500,1 a 1.000 & 8 & 5.432,0 & 4 & 2.934,0 \\
1.000,1 a 2.500 & 8 & 14.974,0 & 0 & --- \\
2.500,1 a 5.000 & 6 & 19.556,0 & 0 & --- \\
5.000,1 a 10.000 & 2 & 12.500,0 & \textit{s} & 5.870,0 \\
\textbf{Más de 10.000} & \textbf{2} & \textbf{37.200,0} & \textbf{0} & \textbf{---} \\
\midrule
\textbf{Total} & \textbf{126} & \textbf{98.774,9} & \textbf{70} & \textbf{11.417,5} \\
\bottomrule
\end{longtable}}

{\small\itshape Para comparar con 2002, el cuadro agrupa en un solo tramo de 5.000,1 a 10.000 hectáreas los dos que el censo de 2018 publica por separado; la explotación de 5.870 hectáreas pertenece al de 5.000,1 a 7.500, como muestra el cuadro anterior.\par}

La pérdida está concentrada en el techo de la escala. En 2002 el departamento tenía \textbf{dos explotaciones de más de diez mil hectáreas}, que sumaban 37.200; dos más de entre cinco y diez mil, con 12.500; seis de entre 2.500 y 5.000, con 19.556; y ocho de entre 1.000 y 2.500, con 14.974. Esos cuatro tramos reunían 84.230 hectáreas: el 85,3 por ciento de la superficie agropecuaria departamental.

En 2018 sólo sobrevive una explotación en esos tramos, la de 5.870 hectáreas. De las 84.230 hectáreas que en 2002 estaban en explotaciones mayores a mil, quedan 5.870.

Los tramos pequeños, en cambio, se mantuvieron o crecieron: las explotaciones de hasta cinco hectáreas pasaron de 18 a 25, y las de cinco a diez, de 4 a 12.

De modo que lo que ocurrió no fue el desplazamiento del pequeño productor por el grande, que es el relato habitual del proceso agrario. Fue lo inverso: \textbf{las grandes explotaciones dejaron de ser explotaciones}, y la minifundización de base persistió. Ochenta mil hectáreas salieron del registro agropecuario.

\begin{cautela}
Qué ocurrió con esas hectáreas es la pregunta que este libro no puede responder, y conviene enumerar las posibilidades sin elegir ninguna.

Pudieron subdividirse y venderse, dejando de constituir unidades productivas --- que es lo que el resto del libro documenta como proceso dominante en el corredor. Pudieron dejar de producir y quedar sin uso, sin cambiar de dueño. Pudieron pasar a régimen de conservación: el departamento contiene el embalse Campo Alegre y áreas categorizadas bajo la Ley 7543 de bosques nativos --- hoy revisada por la Ley 8483, analizada en el capítulo \ref{cap:opacidad} --- y la Ley 7107 de áreas protegidas. Pudieron, simplemente, no ser censadas.
\end{cautela}

Sobre la tercera posibilidad conviene una precisión, porque la Ley 7107 contiene varios instrumentos que este libro no encontró usados en el departamento.

\begin{ficha}{Ley 7107 --- instrumentos disponibles, texto verificado}
\begin{itemize}[leftmargin=*,nosep]
\item \textbf{Art.\ 5 inc.\ e:} entre los objetivos del Sistema, ``conservar los ambientes que circundan \textbf{las nacientes de los cursos de agua}, garantizando su conservación \textbf{a perpetuidad}''.
\item \textbf{Art.\ 5 inc.\ j:} ``minimizar la erosión de suelos, \textbf{garantizando la protección de zonas de alto riesgo de desastres naturales}''.
\item \textbf{Arts.\ 60 y 61:} la declaración procede \textbf{de oficio o a petición de parte}; la Autoridad de Aplicación evalúa, delimita, categoriza y eleva al Poder Ejecutivo, que \textbf{emite el acto administrativo que declara el área}. No requiere ley.
\item \textbf{Art.\ 62:} declarada un área por instrumento del Poder Ejecutivo, \textbf{no puede ser alterada sino por ley}, bajo pena de nulidad.
\item \textbf{Art.\ 26:} las Reservas Naturales Municipales son áreas de dominio municipal, y ``la Provincia podrá realizar convenios con los Municipios, \textbf{para transferir terrenos fiscales} de interés particular para la conservación''.
\item \textbf{Art.\ 47 inc.\ 2:} los propietarios privados que incorporen voluntariamente sus inmuebles pueden acceder a beneficios impositivos provinciales y a \textbf{reducciones de tasas y derechos municipales, previo convenio con las Municipalidades}.
\item \textbf{Art.\ 53:} la adhesión voluntaria es por tiempo indeterminado y la renuncia sólo procede tras veinte años.
\item \textbf{Art.\ 51:} la Autoridad debe impulsar la expropiación cuando las restricciones impidan el ejercicio regular de los derechos.
\item \textbf{Art.\ 52:} Provincia y Municipios pueden declarar un área protegida \textbf{por acta multilateral}, con prohibición de que las signatarias alteren después las condiciones.
\item \textbf{Art.\ 57:} las restricciones y autolimitaciones \textbf{se anotan en los Registros} y siguen al inmueble en toda transferencia. Y ``es deber de la Dirección General de Inmuebles \ldots\ comunicar a la Autoridad de Aplicación la enajenación o transferencia por actos entre vivos del inmueble''.
\end{itemize}
\end{ficha}

\begin{cautela}
Dos de esos artículos describen el departamento sin nombrarlo. El inciso e del artículo 5 manda conservar a perpetuidad los ambientes que circundan las nacientes de los cursos de agua; el inciso j manda proteger las zonas de alto riesgo de desastres naturales. La Caldera es una cuenca de cabecera que se inundó dos veces en trece días, entre diciembre de 2018 y enero de 2019. Existía, desde el año 2000, un sistema legal cuyos objetivos declarados apuntan exactamente ahí.

Y el artículo 60 desarma una idea intuitiva. Crear un área protegida \textbf{no requiere una ley}: basta un acto administrativo del Poder Ejecutivo, a instancia de la Autoridad de Aplicación, de oficio o a pedido de parte. La vía legislativa es una entre varias, y no la más rápida. Una vez declarada por decreto, en cambio, el área sólo puede alterarse por ley --- artículo 62 ---, lo que la vuelve más difícil de deshacer que de hacer.

Que ninguno de esos caminos se haya recorrido en el departamento es el dato. No es que faltara la herramienta ni que la herramienta fuera pesada.
\end{cautela}

\begin{cautela}
Las cuatro explicaciones tienen consecuencias distintas y son todas verificables: la primera con los planos de mensura de la Dirección General de Inmuebles; la segunda y la tercera con el Anexo II de la Ley 8483, que contiene la cartografía vigente del Ordenamiento Territorial de Bosques Nativos; la cuarta con los informes metodológicos del propio censo.

Lo que sí puede afirmarse, y no es poco, es que \textbf{en los mismos años en que el departamento perdió el 88 por ciento de su superficie agropecuaria (2002--2018), casi duplicó su parque habitacional y creció en población más que casi ningún otro distrito del país (2010--2022)}. Las dos series son de fuentes distintas, se miden con metodologías distintas y apuntan en direcciones opuestas sobre el mismo suelo.
\end{cautela}

\begin{cautela}
Existe una hipótesis complementaria sobre el abandono de la producción tabacalera que este libro no puede verificar y que registra por su valor explicativo. El secado de tabaco se realiza en estufas que consumen volúmenes considerables de leña. Si el encarecimiento y la escasez del combustible hubieran vuelto inviable esa etapa del proceso, la consecuencia sería el abandono de tierras en producción por razones estrictamente energéticas, y no por precio del suelo ni por falta de mercado.

La cadena causal resultante sería: ausencia de red de gas $\rightarrow$ dependencia de leña $\rightarrow$ agotamiento y encarecimiento del recurso $\rightarrow$ abandono de la actividad $\rightarrow$ liberación de tierra productiva $\rightarrow$ disponibilidad para loteo.

Si se confirmara, invertiría parcialmente el orden causal que este capítulo propone: no sería el mercado inmobiliario el que desplaza a la producción, sino la carencia de infraestructura energética la que libera la tierra que el mercado después absorbe. Las dos explicaciones no se excluyen y probablemente operen juntas.

\textbf{El apartado siguiente la contrasta con datos censales y la refuta para el período intercensal}: entre 2002 y 2018 la superficie de cultivos industriales del departamento no cayó, creció un 11,7 por ciento. Se conserva el planteo porque el descarte es parcial --- el censo no cubre el período posterior a 2018 --- y porque el modo en que una hipótesis se refuta es información sobre el objeto.
\end{cautela}

\subsection*{Qué se cultivaba, y qué se dejó de cultivar}

Los cuadros de superficie implantada por grupo de cultivos permiten comparar los dos censos.

\needspace{6\baselineskip}
{\sffamily\bfseries\footnotesize Superficie implantada, departamento La Caldera (hectáreas)\par}
\vspace{0.5em}
{\small
\setlength{\LTpre}{0.6em}\setlength{\LTpost}{1.2em}
\begin{longtable}{@{}lrrr@{}}
\toprule
Grupo de cultivos & 2002 & 2018 & Variación \\
\midrule
\endhead
Bosques y montes & 822,5 & 551,0 & $-33,0\%$ \\
\textbf{Cultivos industriales} & \textbf{279,0} & \textbf{311,7} & $\mathbf{+11,7\%}$ \\
Legumbres & 236,8 & --- & $-100\%$ \\
Forrajeras perennes & 223,5 & 82,0 & $-63,3\%$ \\
Forrajeras anuales & 221,0 & 145,9 & $-34,0\%$ \\
Hortalizas & 73,1 & 32,6 & $-55,4\%$ \\
Cereales para grano & 63,6 & 74,0 & $+16,4\%$ \\
Frutales & 21,5 & 14,0 & $-34,9\%$ \\
Aromáticas & 1,5 & --- & $-100\%$ \\
Viveros y flores & 1,2 & 3,6 & --- \\
\midrule
\textbf{Total implantado} & \textbf{1.943,7} & \textbf{1.214,8} & $\mathbf{-37,5\%}$ \\
\bottomrule
\end{longtable}
}

\textit{Fuente: INDEC, CNA 2002 (cuadro 4.4) y CNA 2018, Salta, departamento La Caldera.}

\textbf{Este cuadro refuta la hipótesis energética enunciada más arriba para el período intercensal.}

La hipótesis sostenía que la escasez y el encarecimiento de la leña habían vuelto inviable el secado de tabaco y provocado el abandono de tierras en producción. Si fuera correcta, el renglón de cultivos industriales --- que en el Valle de Lerma corresponde principalmente al tabaco --- debería haberse desplomado entre 2002 y 2018.

Ocurrió lo contrario: creció un 11,7 por ciento, de 279 a 311,7 hectáreas. En un departamento que en el mismo período perdió el 88,4 por ciento de su superficie agropecuaria y el 57,5 de sus explotaciones, la superficie de cultivos industriales aumentó.

La hipótesis queda descartada en los términos en que fue enunciada, y se deja registrada la refutación en lugar de suprimir el planteo.

\begin{cautela}
Dicho eso, corresponde una reconciliación honesta, porque el descarte no es total.

El CNA 2018 releva el período 2017--2018. La observación sobre el abandono de la producción tabacalera por falta de leña es de 2026, y este libro no la tiene en ninguna fuente documental ni periodística: vale sólo como hipótesis. Entre una cosa y otra hay ocho años, y el censo no puede registrar lo que ocurrió después. La hipótesis queda descartada \textbf{para el período intercensal}, no para el presente.

Además, ``cultivos industriales'' es un grupo. Sin un cuadro específico de tabaco no puede establecerse qué proporción de esas 311,7 hectáreas le corresponde. Y una superficie estable es compatible con una fuerte caída en la \textit{cantidad de productores}: menos gente sembrando lo mismo, que es exactamente el proceso que el resto del capítulo documenta.

Verificarlo requiere el cuadro de tabaco por departamento, si existe con ese corte, y los registros de la Cámara del Tabaco de Salta, que lleva padrón de productores por zona.
\end{cautela}

Lo que el cuadro sí muestra, y es un hallazgo por derecho propio, es \textbf{qué} desapareció.

Las legumbres pasaron de 236,8 hectáreas a cero. Las forrajeras perennes cayeron un 63,3 por ciento; las anuales, un 34; las hortalizas, un 55,4; los frutales, un 34,9; los bosques y montes implantados, un 33. Lo único que creció, además de los industriales, fueron los cereales, con 74 hectáreas, y ---en escala mínima--- los viveros y flores.

Es decir: no desapareció la agricultura comercial especializada. Desapareció \textbf{la agricultura mixta}. Lo que se dejó de sembrar fue el poroto, la alfalfa, la verdura y la fruta --- los cultivos de consumo, de rotación y de subsistencia, los que sostienen población rural residente. Lo que quedó en pie fue el cultivo industrial que se vende afuera.

El contraste con la provincia lo subraya: entre 2002 y 2018 Salta aumentó su superficie implantada total un 38,2 por ciento y sus cultivos industriales un 226,7. La Caldera perdió un 37,5 por ciento de superficie implantada mientras sus industriales crecían apenas un 11,7. El departamento no se sumó al ciclo agrícola provincial: se retiró de él, conservando sólo un núcleo.

Y ese retiro tiene un correlato demográfico que el capítulo \ref{cap:poblacion} documenta: para 2010 el departamento registraba la tasa de ocupación más alta y la de inactividad más baja de la provincia, con una población que trabaja fuera. La agricultura que desapareció es la que empleaba gente en el lugar.

\subsection*{El deslinde}

La proporción de explotaciones sin límites definidos también se movió, y en dirección contraria a la pérdida.

{\small
\begin{longtable}{@{}lrrr@{}}
\toprule
 & Total EAP & Sin límites & \% \\
\midrule
\endhead
Salta, CNA 1988 & 9.229 & 4.431 & 48,01 \\
Salta, CNA 2002 & 10.297 & 4.722 & 45,86 \\
Salta, CNA 2018 & 8.705 & 2.851 & 32,75 \\
\midrule
La Caldera, CNA 2002 & 308 & 182 & \textbf{59,09} \\
La Caldera, CNA 2018 & 131 & 61 & \textbf{46,56} \\
\bottomrule
\end{longtable}}

En 2002, el 59,09 por ciento de las explotaciones caldereñas carecía de límites definidos, contra el 45,86 provincial. En 2018 la proporción bajó al 46,56, pero la provincia había bajado al 32,75. El departamento estuvo siempre por encima, y sigue estándolo: hoy registra el nivel que la provincia tenía en 2002.

Y la caída departamental --- de 182 explotaciones sin límites a 61 --- fue más pronunciada que la de las explotaciones con límites, que pasaron de 126 a 70. Es decir que la desaparición de unidades productivas golpeó \textbf{más fuerte sobre la tenencia sin deslinde}, que es la más precaria.

Sesenta y una explotaciones sin límites definidos siguen existiendo en el departamento. Esa población no aparece en ninguna cobertura periodística relevada ni en ningún expediente judicial identificado, y la literatura académica la roza sólo a través de la comunidad Kondorwaira (véase más arriba). El Censo Agropecuario la cuenta, y registra entre los dos censos un saldo de ciento veintiuna menos: de 182 a 61.

\subsection*{En 1980 había tabaco, y el Estado lo dijo por decreto}

Este capítulo discute, con dificultad, cuánto tabaco se cultivaba en el departamento. La
hipótesis de que su colapso explique la pérdida de superficie agropecuaria quedó refutada por
el dato censal para el período intercensal, y acotada por el mapa tabacalero provincial. Un decreto de 1980 aporta la
prueba directa que faltaba.

\begin{ficha}{Decreto Nº 12 del 4 de enero de 1980}
Expediente C/24-02529/79. Visto que la \textbf{Dirección General Agropecuaria} eleva
evaluación de \textbf{daños por granizo} en las localidades de \textbf{La Caldera y Chicoana},
y a fin de incluir «el \textbf{área tabacalera} de los Departamentos afectados» en el régimen
de \textbf{Emergencia Agropecuaria}, se declara zona de emergencia por \textbf{seis meses}
desde el 1º de diciembre de 1979.

Los productores tienen \textbf{sesenta días hábiles} para denunciar, y el umbral de daño
exigido es de \textbf{más del cincuenta por ciento}.

Firma el gobernador Ulloa.
\end{ficha}

\textbf{Primera: en 1979 había área tabacalera en La Caldera, y el Estado la reconocía como
tal.} No es una inferencia sobre superficies implantadas: es un decreto que la nombra para
protegerla. \textbf{El departamento aparece junto a Chicoana}, que es uno de los distritos
tabacaleros centrales de la provincia.

\textbf{Segunda: había productores organizados como para que el régimen tuviera sentido.} Un
decreto de emergencia con plazo de denuncia y umbral de daño supone gente que va a denunciar.
El capítulo documenta que en 2015 convocaba a asamblea la Asociación Civil de Pequeños
Productores Agropecuarios del departamento ---y siguió haciéndolo al menos hasta 2024---; en 1979 había a quién indemnizar.

\textbf{Y tercera, que es la que importa para la tesis de este capítulo: eso no cambia la
refutación, la precisa.} Que hubiera tabaco en 1979 no significa que su desaparición explique
la pérdida del ochenta y ocho por ciento de la superficie agropecuaria entre 2002 y 2018,
porque en ese período la superficie de cultivos industriales del departamento \textit{creció}.
\textbf{Lo que el decreto agrega es el punto de partida de una serie que el libro no puede
reconstruir}: había área tabacalera reconocida en 1979, y en 2018 el censo sólo registra
cultivos industriales sin desagregar, de modo que no permite aislarla; los cuarenta
años intermedios están en el hueco que el capítulo \ref{cap:metodo} declara.

\section{Qué era esa superficie: el censo de 2008}

Este capítulo compara dos totales ---98.774 hectáreas en 2002 y 11.417 en 2018--- y de esa
comparación saca una caída del ochenta y ocho por ciento. \textbf{El censo intermedio permite
ver de qué estaba hecha esa superficie, y el resultado obliga a matizar la lectura.}

\begin{ficha}{Censo Nacional Agropecuario 2008, provincia de Salta --- departamento La Caldera}
\textbf{Explotaciones agropecuarias: 174}, de las cuales \textbf{86 sin límites definidos} y
\textbf{88 con límites definidos}. Superficie de estas últimas: \textbf{75.870,6 hectáreas}.

\textbf{Por escala de extensión} ---sobre las 88 con límites---: catorce de hasta 5 ha; trece de
5 a 10; diecinueve de 10 a 25; dos de 25 a 50; cinco de 50 a 100; cinco de 100 a 200; diez de
200 a 500; seis de 500 a 1.000; cinco de 1.000 a 2.500; seis de 2.500 a 5.000; dos de 5.000 a
10.000; y \textbf{una de más de 10.000 hectáreas}.

\textbf{Por uso de la tierra:} superficie \textbf{implantada 2.018 ha ---el 2,7 por ciento---};
superficie destinada a otros usos \textbf{73.852,6 ---el 97,3 por ciento---}.

Dentro de lo implantado: cultivos anuales 314,1; cultivos perennes 34; forrajeras anuales 348;
forrajeras perennes 207,2; bosques y montes implantados 1.060; sin discriminar 54,7.

Dentro de los otros usos: \textbf{bosques y montes espontáneos 39.613,3 ---el 52,2 por ciento de la superficie de las explotaciones con límites---}; \textbf{pastizales 15.136 ---19,9 por ciento---}; sin
discriminar 8.496,1; no apta o de desperdicio 6.377,3; apta no utilizada 3.944,1; caminos,
parques y viviendas 285,8.
\end{ficha}

\begin{cautela}
\textbf{Esto cambia el sentido de la cifra que abre el capítulo, y conviene decirlo sin
rodeos.}

\textbf{La «superficie agropecuaria» de este departamento era, en su abrumadora mayoría, monte y
pastizal.} Sólo el 2,7 por ciento estaba implantado con algo, y dentro de ese 2,7 por ciento la
mitad son bosques y montes implantados: \textbf{los cultivos anuales y perennes juntos no llegaban a
cuatrocientas hectáreas en todo el departamento}, sobre setenta y cinco mil declaradas.

\textbf{De modo que la caída del ochenta y ocho por ciento no mide la desaparición de una
agricultura.} Mide, sobre todo, que \textbf{dejaron de declararse como explotación
agropecuaria grandes extensiones de monte y ladera}. Una sola unidad de más de diez mil
hectáreas explica por sí misma una porción importante del total de 2008.

\textbf{Y eso vuelve consistente lo que documenta el capítulo \ref{cap:trabajo}}: que después de
la caída todavía hay ciento noventa y tres personas ocupadas en agricultura. Si lo que se
perdió fue superficie declarada de monte y no superficie cultivada, no había por qué esperar un
derrumbe del empleo agrícola. \textbf{Las dos cosas eran compatibles y este libro no lo había
visto.}
\end{cautela}

\begin{cautela}
\textbf{Tres límites de este cuadro.}

\textbf{El CNA 2008 tuvo una cobertura deficiente reconocida}, y el INDEC nunca publicó
resultados definitivos de ese operativo. Sus cifras deben leerse como indicativas de
composición, no como magnitudes exactas.

\textbf{Ochenta y seis de las ciento setenta y cuatro explotaciones no tenían límites
definidos}, y su superficie no se computa. Es la mitad de las unidades del departamento, y la
mitad sobre la que menos se sabe: son, típicamente, las más chicas y las de tenencia precaria.

\textbf{Y este libro no sabe qué pasó con esas 75.870 hectáreas.} Dejar de declararse como
explotación agropecuaria no es lo mismo que cambiar de uso ni de dueño. El monte puede seguir
siendo monte y el titular el mismo. \textbf{La caída del censo es un hecho registral antes que
territorial}, y el apéndice \ref{cap:pedidos} pide lo que permitiría distinguirlos.
\end{cautela}

\subsection*{Qué peso tiene el tabaco en este departamento}

La hipótesis enunciada más arriba depende de que el tabaco haya sido una porción significativa de la superficie agrícola del departamento. Ese supuesto no pudo confirmarse, y lo que sí se obtuvo apunta en contra.

\begin{ficha}{La Caldera en el mapa tabacalero provincial}
\begin{itemize}[leftmargin=*,nosep]
\item Un estudio del Ministerio de Agricultura de la Nación elaborado sobre el \textbf{Censo Nacional Agropecuario 2002} ubica las explotaciones tabacaleras salteñas concentradas en el valle de Lerma y, ``en menor medida'', en La Caldera, La Candelaria, La Viña, Metán y Guachipas.
\item Un parte del gobierno provincial de 2023, con datos de la Cámara del Tabaco de Salta, consigna que el \textbf{90 por ciento} de la producción se distribuye en General Güemes, Rosario de Lerma, Chicoana y Cerrillos, y que el \textbf{10 por ciento restante} se reparte entre Metán, Guachipas, Capital, \textbf{La Caldera}, La Viña y La Candelaria.
\end{itemize}
\end{ficha}

\begin{cautela}
En las dos puntas del período que este capítulo mide --- 2002 y el presente --- La Caldera aparece del mismo lado: en la cola del mapa tabacalero, compartiendo con otros cinco departamentos una décima parte de la producción provincial.

Eso vuelve improbable que la retracción del tabaco explique el desplome del 88,4 por ciento de la superficie agropecuaria departamental ---que es, sobre todo, monte declarado---, y ni siquiera la caída del 37,5 por ciento de la superficie implantada. Un cultivo que nunca fue dominante en un territorio no puede, al retirarse, arrastrar la mayor parte de lo que se cultivaba.

\textbf{Pero improbable no es imposible, y conviene decir por qué.} Los dos datos disponibles miden participación en la producción \textbf{provincial}, no participación del tabaco en la superficie \textbf{del propio departamento}. Un departamento con poca superficie agrícola puede tener en ella una proporción alta de tabaco y aun así pesar poco en el total de Salta. Las dos cosas son compatibles y sólo el cuadro de superficie implantada por cultivo a nivel departamental permite distinguirlas.

Este libro, entonces, la tiene por refutada para 2002--2018 y no la descarta para el período posterior: la deja en suspenso con una presunción en contra, y señala exactamente qué dato la resolvería. Si al obtenerlo el tabaco resultara marginal también dentro del departamento, la explicación del desplome habrá que buscarla en la combinación que el censo de 2008 sugiere: grandes extensiones de monte que dejaron de declararse como explotación, y una parte convertida en suelo urbano, que es lo que el resto de este libro documenta.
\end{cautela}

\section{Hubo un aserradero en el pueblo}

Este capítulo discute, entre otras cosas, la hipótesis de que la escasez de leña haya
incidido en la retracción de la actividad agropecuaria. El archivo aporta un dato sobre el
otro extremo de esa cadena: la explotación forestal existía y tenía escala industrial.

\begin{ficha}{Remate judicial sin base --- B.O.\ Nº 1.650, 21 de agosto de 1936}
Por disposición del Juez de Comercio, Dr.\ Néstor Cornejo Isasmendi, en el juicio ejecutivo
seguido por \textbf{Francisco Moschetti y Compañía contra don Rafael Delgado\index{Delgado, Rafael}}, el martillero
Ernesto Campilongo venderá \textbf{sin base y de contado} los bienes que se encuentran en
poder del ejecutado \textbf{en el pueblo La Caldera}:

«Una \textbf{sierra sin fin}; dos \textbf{circulares}; un \textbf{motor a vapor}; una
\textbf{garlopa}; un lote de \textbf{orilleros desperdicios} y \textbf{diez y seis bueyes}.»
\end{ficha}

El inventario describe un \textbf{aserradero completo}: la sierra principal, dos sierras
circulares para escuadrar, la fuente de energía, una máquina de cepillado, el descarte de
orillas ya cortado y la yunta de dieciséis bueyes para el arrastre de rollizos.

Tres cosas se siguen de ahí.

\textbf{Había industria forestal en el pueblo, no sólo extracción de leña.} Un motor a vapor
y una garlopa no sirven para hacer carbón: sirven para producir madera aserrada y cepillada.

\textbf{¿Y el ejecutado era el juez de paz?} En las mismas ediciones de 1936, el nombramiento
de jueces de paz del distrito municipal de La Caldera consigna a un \textbf{Rafael Delgado\index{Delgado, Rafael}}.
Es el mismo nombre del dueño del aserradero cuyos bienes se remataron sin base ese año.

\textbf{Y hay un tercer punto, diecinueve años anterior.} En septiembre de 1917 el decreto 1377 nombra miembro de la comisión municipal de La Caldera, en reemplazo de Andrés Marini que renunció, a un \textbf{Rafael Delgado}\index{Delgado, Rafael} (B.O.\ Nº 668, h.\ 2; capítulo \ref{cap:siglo}). Tres apariciones del mismo nombre en el mismo departamento: comisión municipal en 1917, juez de paz en 1936 y dueño del aserradero ejecutado ese mismo año.

Como siempre, coincidencia de nombre y no identificación registral. Pero si se tratara de la
misma persona --- y el apéndice \ref{cap:personas} reúne otras trayectorias del mismo tipo --- el juez de
paz del pueblo habría perdido su aserradero en subasta pública a instancia de una casa
comercial de Salta, en el mismo año en que ejercía el cargo.

\textbf{Y se remató sin base, por deuda con una casa comercial de la capital.} Como la
fundición de Mojotoro en 1925, como las dos casas quinta, como Potrero de Gallinato y como La
Despensa. La actividad económica del departamento aparece en el registro público,
sistemáticamente, en el momento en que se liquida.

\textbf{Con esto, el arco económico del departamento en la primera mitad del siglo queda así:}
una fundición metalúrgica de plomo entre 1917 y 1925, ocho solicitudes de cateo minero en
1925 que nadie exploró, una única mina registrada que en 1933 seguía sin mensurar, un
aserradero rematado en 1936, y una incorporación a la región económica del olivo en 1934 ---que el apartado siguiente presenta--- que
el Censo Agropecuario no registra entre lo que hoy se cultiva.

\textbf{Cinco actividades, cinco intentos, ninguno en pie.} La conversión del suelo a uso
residencial no desplazó a una economía productiva próspera: llegó sobre el sedimento de una
serie de actividades que habían llegado de afuera, durado poco y liquidado en subasta.

\section{Qué le pasó a la gente que trabajaba esa tierra}

Este capítulo mide hectáreas. Conviene cerrar preguntando por las personas, porque el
archivo permite hoy formular la pregunta mejor que cuando este trabajo empezó.

En 2015 existía ---y en 2024 seguía convocando asambleas--- una asociación civil de pequeños productores del departamento, con personería y sede en Vaqueros (capítulo \ref{cap:resistencias}).

El desplome del ochenta y ocho por ciento de la superficie agropecuaria no ocurrió sobre un
vacío: ocurrió sobre productores organizados. Y ocho décadas antes, en 1934, la provincia
había incorporado a La Caldera a la «región económica del olivo» y eximido por diez años de
todo impuesto a las tierras donde se plantaran olivos --- un intento de asignarle una
vocación productiva que el Censo Agropecuario no registra entre lo que hoy se cultiva.

\textbf{La pregunta de este capítulo deja de ser qué pasó con las hectáreas y pasa a ser qué
pasó con la gente que las trabajaba.} Este libro no la contesta: no hizo entrevistas sobre eso, no
consultó el registro de personas jurídicas ni los padrones del Censo Agropecuario a nivel de
explotación. Pero ahora sabe dónde empezar a buscar, y eso está en el apéndice \ref{cap:pedidos}.

\section{Los pobladores que la periodización omite}

La literatura académica identifica una secuencia de movimientos poblacionales que desaparece de toda la cobertura periodística posterior: habitantes originarios; primeras familias de origen europeo instaladas como pequeños agricultores; migración desde el interior provincial y desde Bolivia de trabajadores del tabaco, primero golondrina y luego estabilizada; y desde los años noventa, la migración urbano-periurbana de sectores medios que acompaña el loteo.

Los radios censales correspondientes a Vaqueros muestran población integrada por trabajadores del tabaco que en los años sesenta accedieron a loteos de vivienda unifamiliar, y por antiguos pobladores que habían accedido a las tierras en calidad de piseros. Sobre esa base social se superpone, treinta años después, un mercado orientado a sectores medios y medios-altos.

El resultado que la literatura consigna es un desplazamiento en sentido inverso: la subdivisión, el loteo y el aumento del precio de la tierra, sumados a la caída de la demanda de mano de obra local, expulsan población desde Vaqueros hacia los asentamientos periféricos de la capital.

Ésta es la pieza que ninguna crónica sobre el ``boom inmobiliario'' menciona y que el libro sostiene como eje: el mismo proceso que trae vecinos nuevos expulsa vecinos viejos. El capítulo \ref{cap:poblacion} muestra el resultado en datos censales: para 2010 el departamento ya tenía la tasa de ocupación más alta y la de inactividad más baja de la provincia, con hogares más chicos que el promedio. La población de los piseros y los trabajadores del tabaco había sido reemplazada por otra que trabaja en la capital y duerme en el corredor. No son dos fenómenos, es uno. El capítulo \ref{cap:tierrafiscal} mostrará que el circuito estatal de tierra fiscal, lejos de compensar ese desplazamiento, en ocasiones lo formaliza.
